\documentclass[a4paper,12pt,answers]{exam}
\usepackage{graphicx} % Required for inserting images

\usepackage{amsmath}
\usepackage{amssymb}

\title{CS215 Assignment-1}
\author{Malay Kedia (Roll No. 23B1034) \\ Aakash Gupta (Roll No. 23B0953) \\ Yash Singh (Roll No. 23B0934)}
\date{}

\begin{document}
\maketitle

\begin{questions}
\question{\bf A Pretty “Normal" Mixture}

\textbf{Task A:}
\begin{solution}\\
\begin{align*}
F_{A}(c)&=P(A<c)\\
&=\sum_{i=1}^{K} P(A < c \land A = X_i) \hspace{3cm} \text{(By law of total probability)}\\
&=\sum_{i=1}^{K} P(A < c \mid A = X_i) \cdot P(A = X_i)\\
&=\sum_{i=1}^{K} P(X_i < c) \cdot P(A = X_i)\\
&= \sum_{i=1}^{K} p_i F_{X_i}(c)
\end{align*}
(Substituting the CDF $F_{X_i}(c)$ for each Gaussian component)

To find the probability density function (PDF) $f_A(c)$, we differentiate $F_A(c)$ with respect to $c$.
\begin{align*}
f_A(c) &= \frac{d}{dc} F_A(c)\\
&= \sum_{i=1}^{K} p_i \cdot \frac{d}{dc} F_{X_i}(c)\\
&= \sum_{i=1}^{K} p_i \cdot f_{X_i}(c) \hspace{3cm} (\frac{d}{dc} F_{X_i}(c) = f_{X_i}(c))
\end{align*}
The PDF of the GMM variable X is $f_x$, given by $f_X(x)=\sum_{i=1}^{K} p_i \cdot f_{X_i}(x)$, which is exactly same as the one I got from my algorithm. Hence, $\forall u \in R, f_A(u)=f_X(u)$.
\end{solution}

\textbf{Task B:}
\begin{solution}\\
Part 3:\\
\begin{align*}
M_X(t)&=E(e^{tX})\\
&=\int_{-\infty}^{\infty} e^{tx} f_X(x) dx\\
&=\int_{-\infty}^{\infty} e^{tx}  \cdot (\sum_{i=1}^{K} p_i \cdot f_{X_i}(x)) dx\\
&=\sum_{i=1}^{K} (p_i \int_{-\infty}^{\infty} e^{tx} \cdot f_{X_i}(x) dx )\\
&=\sum_{i=1}^{K} p_i E(e^{t X_i})\\
&=\sum_{i=1}^{K} p_i \cdot e^{\mu_i t + \frac{\sigma_i^2 t^2}{2}}
\end{align*}

We can now find the expected value $E[X]$ and variance $Var[X]$ using the MGF.

Part 1:\\
Differentiating the MGF with respect to t:

\begin{align*}
M'_X(t) &= \frac{d}{dt} \left( \sum_{i=1}^{K} p_i \cdot e^{\mu_i t + \frac{\sigma_i^2 t^2}{2}} \right)  \\
&= \sum_{i=1}^{K} p_i \cdot \frac{d}{dt} \left( e^{\mu_i t + \frac{\sigma_i^2 t^2}{2}} \right)\\
&=\sum_{i=1}^{K} p_i \cdot e^{\mu_i t + \frac{\sigma_i^2 t^2}{2}} \cdot \left( \mu_i + \sigma_i^2 t \right)\\
\end{align*}


Now, $E[X]=M'_X(0)$

\begin{align*}
E[X] &= M'_X(0) = \sum_{i=1}^{K} p_i \cdot e^{0} \cdot \mu_i \\
&= \sum_{i=1}^{K} p_i \cdot \mu_i
\end{align*}

Part 2:
\begin{align*}
M''_X(t) &= \sum_{i=1}^{K} p_i \cdot \frac{d^2}{dt^2} \left( e^{\mu_i t + \frac{\sigma_i^2 t^2}{2}} \right)\\
&= \sum_{i=1}^{K} p_i \cdot e^{\mu_i t + \frac{\sigma_i^2 t^2}{2}} \cdot \left( \mu_i + \sigma_i^2 t \right)^2 + \sum_{i=1}^{K} p_i \cdot e^{\mu_i t + \frac{\sigma_i^2 t^2}{2}} \cdot \sigma_i^2
\end{align*}

Now, the second moment $E[X^2]$ is:
\[
E[X^2] = M''_X(0) = \sum_{i=1}^{K} p_i \cdot \left( \mu_i^2 + \sigma_i^2 \right)
\]

Now, the variance $\text{Var}(X)$ is:
\[
\text{Var}(X) = E[X^2] - (E[X])^2
\]
\[
\text{Var}(X) = \sum_{i=1}^{K} p_i \cdot \left( \mu_i^2 + \sigma_i^2 \right) - \left( \sum_{i=1}^{K} p_i \cdot \mu_i \right)^2
\]
\end{solution}

\textbf{Task C:}
\begin{solution}\\
ATQ, $Z=\sum_{i=1}^K p_i X_i$.

Part 4:
\begin{align*}
M_Z(t)&=E[e^{tZ}]\\
&=E[e^{t\sum_{i=1}^K p_i X_i}]\\
&=E[\prod_{i=1}^K e^{t p_i X_i}]\\
&=\prod_{i=1}^K E[e^{t p_i X_i}] \hspace{3cm} \text{(As variables are independant)}\\
&=\prod_{i=1}^K e^{\mu_i p_i t + \frac{\sigma_i^2 p_i^2 t^2}{2}}\\
&=e^{\sum_{i=1}^K \mu_i p_i t + \frac{\sigma_i^2 p_i^2 t^2}{2}}\\
&=e^{ t \sum_{i=1}^K \mu_i p_i + \frac{t^2 \sum_{i=1}^K \sigma_i^2 p_i^2}{2}}\\
\end{align*}


Part 1:\\
Differentiating $M_Z(t)$ with respect to $t$:

\begin{align*}
M'_Z(t) &= \frac{d}{dt} \left(e^{t \sum_{i=1}^K \mu_i p_i + \frac{t^2 \sum_{i=1}^K \sigma_i^2 p_i^2}{2}}\right)\\
&= e^{t \sum_{i=1}^K \mu_i p_i + \frac{t^2 \sum_{i=1}^K \sigma_i^2 p_i^2}{2}} \cdot \left(\sum_{i=1}^K \mu_i p_i + t \sum_{i=1}^K \sigma_i^2 p_i^2\right)    
\end{align*}

\begin{align*}
E[Z] &= M'_Z(0) = e^0 \cdot \left(\sum_{i=1}^K \mu_i p_i\right) \\
&= \sum_{i=1}^K \mu_i p_i
\end{align*}

Part 2:
\begin{align*}
M''_Z(t) &= \frac{d}{dt} \left( e^{t \sum_{i=1}^K \mu_i p_i + \frac{t^2 \sum_{i=1}^K \sigma_i^2 p_i^2}{2}} \cdot \left(\sum_{i=1}^K \mu_i p_i + t \sum_{i=1}^K \sigma_i^2 p_i^2\right) \right)\\
&= e^{t \sum_{i=1}^K \mu_i p_i + \frac{t^2 \sum_{i=1}^K \sigma_i^2 p_i^2}{2}} \cdot \left( \left(\sum_{i=1}^K \mu_i p_i + t \sum_{i=1}^K \sigma_i^2 p_i^2\right)^2 + \sum_{i=1}^K \sigma_i^2 p_i^2 \right)
\end{align*}

Evaluating at 0,
\[
E[Z^2]= M''_Z(0) = e^0 \cdot \left( \left(\sum_{i=1}^K \mu_i p_i\right)^2 + \sum_{i=1}^K \sigma_i^2 p_i^2 \right) = \left(\sum_{i=1}^K \mu_i p_i\right)^2 + \sum_{i=1}^K \sigma_i^2 p_i^2
\]

Now, the variance is:
\begin{align*}
\text{Var}(Z) &= E[Z^2] - (E[Z])^2\\
&= \left(\sum_{i=1}^K \mu_i p_i\right)^2 + \sum_{i=1}^K \sigma_i^2 p_i^2 - \left(\sum_{i=1}^K \mu_i p_i\right)^2\\
&= \sum_{i=1}^K \sigma_i^2 p_i^2
\end{align*}


Part 3: PDF of $Z$

We observe that the MGF of Z matches the MGF of the gaussian $\mathcal{N}(\mu, \sigma^2)$ with $\mu=\sum_{i=1}^K \mu_i p_i$ and $\sigma^2=\sum_{i=1}^K \sigma_i^2 p_i^2$.

\[
Z \sim \mathcal{N}\left(\sum_{i=1}^K \mu_i p_i, \sum_{i=1}^K \sigma_i^2 p_i^2\right)
\]

Thus, the PDF of $Z$ is:

\[
f_Z(u) = \frac{1}{\sqrt{2\pi \sum_{i=1}^K \sigma_i^2 p_i^2}} \exp\left(-\frac{(u - \sum_{i=1}^K \mu_i p_i)^2}{2 \sum_{i=1}^K \sigma_i^2 p_i^2}\right)
\]

Part 5: \\
No, X and Z have very different behaviors, and dont even have the same variances, much less the PDFs.

Part 6:\\
As concluded in part 3, Z is a gaussian.
\end{solution}

\textbf{Task D:}
\begin{solution}\\
Let $X$ be a finite discrete random variable taking values $x_1, x_2, \dots, x_n$ with respective probabilities $P(X = x_i) = p_i$, for $i = 1, 2, \dots, n$.\\

% Consider any random variable Z. Its moment generating function (MGF) is given by  
% \[
% M_Z(t) = \sum_{k=0}^{\infty} \frac{E[Z^k]}{k!} t^k,
% \]
% This can be proved as follows:

% \begin{align*}
% M_Z(t) &= E[e^{tZ}]\\
% &= E\left[\sum_{k=0}^{\infty} \frac{(tZ)^k}{k!}\right] \hspace{2cm} \text{(As $e^{tZ} = \sum_{k=0}^{\infty} \frac{(tZ)^k}{k!}$)}\\
% &= \sum_{k=0}^{\infty} E\left[\frac{(tZ)^k}{k!}\right] \hspace{2cm} \text{(By Linearity of Expectation)}\\
% &= \sum_{k=0}^{\infty} \frac{t^k}{k!} E[Z^k]
% \end{align*}

The moment generating function (MGF) of $X$, denoted by $\phi_X(t)$, is given by:
\[
\phi_X(t) = E[e^{tX}] = \sum_{i=1}^n e^{t x_i} P(X = x_i) = \sum_{i=1}^n e^{t x_i} p_i
\]

By the above construction, it is trivial to show that if PMF of two random variables is same, their MGF is the same. Hence, the PMF uniquely determines the MGF.

% Consider the MGF $\phi_X(t)$. Since the MGF is a power series in $t$ (the sum of exponentials of the discrete values $x_i$), it can be expanded as a Taylor series around $t = 0$ as done above. The coefficients of this series are the moments of the distribution, which are directly related to the probabilities $p_1, p_2, \dots, p_n$.
% \[
% E[Z^k]=\sum_{i=1}^n p_i x_i^k
% \]

Given the MGF $\phi_X(t) = \sum_{i=1}^n e^{t x_i} p_i$, we can determine the probability mass function (PMF) by solving a system of linear equations. If we know $\phi_X(t)$ for enough values of $t$, we can solve for the probabilities $p_1, p_2, \dots, p_n$ by inverting this system. Therefore, the MGF uniquely determines the PMF.

Conclusion about $X$ and $Z$

Since $Z$ is a weighted sum of independent Gaussian random variables and does not have the same MGF as $X$, they dont have the same distribution (by the theorem). This is also the expected logical outcome, as $X$ is a GMM, in which probability of the rv itself is a weighted sum of probabilities of n gaussian random variables, while $Z$ is the sum of outcomes of n independant gaussian random variables, which turns out of itself be a gaussian.
\end{solution}

\end{questions}

\end{document}
